\subsection{$(B,A,A)$ 型膜}

对 $(B,B,B)$ 型膜施加 $S$-对偶，便得到 $(B,A,A)$ 型膜。一般说来，
如何把某个特定的膜在 $S$-对偶下作变换可能并不清楚，但把未破缺的
超对称在 $S$-对偶下作变换却毫无困难。由式 \eqref{eq:tokio} 可见，
在 $S$ 变换下（取 ${\rm Re}\,\tau=0$），超对称的生成元 $\epsilon_{\ell,r}$
变换为 % 式 (12.20)
\setcounter{equation}{19}
\begin{equation}\label{eq:bilj}
\epsilon_\ell\to
\epsilon_\ell\cdot\exp\left(-{i\pi\over 4}\right),\quad
\epsilon_r\to \epsilon_r\cdot\exp\left({i\pi\over 4}\right),
\end{equation}
 或者等价地 % 式 (12.21)
\setcounter{equation}{20}
\begin{equation}\label{eq:yogo}
\epsilon\to\epsilon'={1+\HG\over \sqrt
2}\epsilon
\end{equation}
 其中 $\HG=\Gamma_0\Gamma_1\Gamma_2\Gamma_3$ 度量四维手征性。若
$T\epsilon=\epsilon$，则
$\epsilon'={1+\HG\over \sqrt 2}\epsilon={1+\HG\over \sqrt
2}T\epsilon={1+\HG\over \sqrt 2}T{1-\HG\over \sqrt 2}\epsilon'$，
于是 $T'\epsilon'=\epsilon'$，其中 % 式 (12.22)
\setcounter{equation}{21}
\begin{equation}\label{eq:uvuf}
T'={1+\HG\over \sqrt
2}T{1-\HG\over \sqrt 2}.
\end{equation}
 若 $T=\Gamma_0\Gamma_4$，则
% 式 (12.23)
\setcounter{equation}{22}
\begin{equation}\label{eq:polyo}
T'=\Gamma_1\Gamma_2\Gamma_3\Gamma_4.
\end{equation}

我们在费米子上施加边界条件 $T'\lambda|=-\lambda|$。于是超对称要求
$\{T',\Gamma^{IJ}F_{IJ}\}\epsilon=0$，即
% 式 (12.24)
\setcounter{equation}{23}
\begin{equation}\label{eq:bngo}
\left(\sum_{I,J\in
\{1,2,3,4\}}\Gamma^{IJ}F_{IJ}+\sum_{I,J\notin
\{1,2,3,4\}}\G^{IJ}F_{IJ}\right)\epsilon=0.
\end{equation}

这一条件意味着对 $I>4$ 有 $F_{0I}|=0$，因此我们对 $I>4$ 的 $A_I$
施加 Dirichlet 边界条件。我们干脆令 $A_I\vert=0$（$I=5,8,9$），
并寻找使「Higgs 场」$\phi_2=A_6,$ $\phi_3=A_7$ 取非零值的边界
条件。假定 \eqref{eq:bngo} 应对所有满足
$\Gamma_{1234}\epsilon=\epsilon$ 的 $\epsilon$ 成立（不再附带
$\G_{2367}\epsilon=\epsilon$ 这一条件），则该式意味着
$[\phi_2,\phi_3]=F_{67}=0$。在此基础上，$\phi_2$ 与 $\phi_3$ 的
规范不变内容体现在 $\phi_z={1\over 2}(\phi_2-i\phi_3)$ 的特征
多项式中。我们选取 Dirichlet 边界条件，规定该特征多项式。

另一方面，$A_{2}$, $A_{3}$, $A_{4}$ 满足修正的 Neumann 边界条件 % 式 (12.25)
\setcounter{equation}{24}
\begin{equation}\label{eq:yimgo}
\begin{aligned} F_{12}\vert -
F_{34}\vert
 &=0\\
 F_{13}\vert - F_{42}\vert
 &=0\\
 F_{14}\vert - F_{23}\vert
 &=0,\\\end{aligned}
\end{equation}

它们来自对 \eqref{eq:bngo} 施加 $\Gamma_{1234}\epsilon=\epsilon$。

我们还要求满足 $\Gamma_{2367}\eps=\eps$ 的超对称不被破坏。这给
满足 Dirichlet 条件的场的边界值加上限制。要求费米子的超对称
变分在边界上为零，得到
% 式 (12.26)
\setcounter{equation}{25}
\begin{equation}\label{eq:dirbaa}
\begin{aligned}F_{26}|-F_{73}|&=0,\\ F_{27}|-F_{36}|&=0.\end{aligned}
\end{equation}

这些方程等价于 $D_\bz\phi_z=0$，即「复」Hitchin 方程。

这些边界条件意味着：在由 $(A_{\bar z},\phi_z)$ 对确定的 Higgs 丛
$E$ 中，决定 $E$ 全纯结构的场 $A_{\bar z}$ 可以涨落，而 $\phi_z$ 的
特征多项式保持固定。很可能相应的膜支撑在 Hitchin 映射
$\pi:\MH\to {\EUBB}$ 的一个纤维上。这给了我们一个原型的
$(B,A,A)$ 型膜。

这一经典分析中出现的条件 $[\phi_2,\phi_3]=0$ 并不是定义 $\MH$
的 Hitchin 方程之一。我们本应有 $[\phi_2,\phi_3]-F_{23}=0$。
很可能这一额外项是在向红外重整化群流的过程中作为某种量子修正
出现的。

在规范场论退化为二维 sigma 模型的极限下，$S$-对偶把缠绕 Hitchin
纤维化纤维的 $(B,A,A)$ 型膜映为 $(B,B,B)$ 型零膜。在阿贝尔情形，
可以直接验证我们为这些膜描述的边界条件互为 $S$-对偶。

我们还可以把 $T'$ 换成 $T''=T'\G_{2367}=-\G_{1467}$，构造另一种
看似同样支撑在 Hitchin 纤维化纤维上的 $(B,A,A)$ 型膜。施加
$T''\epsilon=\epsilon$ 与 $T''\lambda\vert=\lambda\vert$，我们得到
% 式 (12.27)
\setcounter{equation}{26}
\begin{equation}\label{eq:loppo}
[\G^{IJ}F_{IJ},T'']\epsilon=0,
\end{equation}
 这再次意味着对 $A_6=\phi_2$ 与 $A_7=\phi_3$ 的 Dirichlet 边界
条件，以及对 $A_2$ 与 $A_3$ 的 Neumann 边界条件。

当然，对 Hitchin 纤维化 $\pi:\MH\to{{\cmmib B}}$ 的给定纤维
${\CMF}$，应当有许多膜，它们与 ${\CMF}$ 上可能的平坦线丛相联系。
我们给出了两种构造，但由于时间反演对称性，它们各自给出的都是
平凡的 Chan-Paton 线丛（至多是二阶的）。我们希望能修改规范场论
的构造，使该丛可以变动。

\subsection{$(A,B,A)$ 型膜}\label{sec:xgoto}

下一个目标是从规范场论构造一些 $(A,B,A)$ 型膜。其中最令人关注
的是典范余各向同性膜。作为准备，我们引入缩写
$\Gamma_{\pm I\pm J}=(\pm \G_I\pm \G_J)/\sqrt 2$。类似地，取
$F_{\pm I\pm J,K}=(\pm F_{IK}\pm F_{JK})/\sqrt 2$，等等。

让我们定义 % 式 (12.28)
\setcounter{equation}{27}
\begin{equation}\label{eq:jokko}
T'=-\Gamma_{15}\Gamma_{2-7}\Gamma_{3+6}.
\end{equation}

稍后就会看到，这一选择导出 $(A,B,A)$ 型膜。这个定义的一个要点
是：$T'$ 在 $x^2-x^3$ 平面与 $x^6-x^7$ 平面同时旋转下不变。因此在
扭化理论中，只要坐标彼此正交，$T'$ 就不依赖 $x^2,x^3$ 的具体
选取。同理，若在 $M=\RR_+^2\times C$ 上考虑扭化理论，$T'$ 只依赖
$C$ 上的度规，而不依赖具体的坐标选择。我们用 $T'$ 构造的膜只
依赖度规的共形类，即 $C$ 上的复结构。

我们仍可构造两类膜，其边界条件分别为 $T'\lambda|=\lambda|$ 与
$T'\lambda|=-\lambda|$。在前一情形，超对称按通常方式给出
$[\G^{IJ}F_{IJ},T']\epsilon=0$。这告诉我们：若 $I,J$ 之一落在由
$1,5,2-7,$ 与 $3+6$ 方向生成的子空间中，另一个与该子空间正交，
则 $F_{IJ}=0$。特别地，对 $I=5,2-7,3+6$ 有 $F_{0I}|=0$，我们对
这些 $I$ 施加 Dirichlet 边界条件。而对 $I=2+7,3-6,4,8,9$ 有
$F_{1I}=0$，我们对这些 $I$ 的 $A_I$ 施加 Neumann 边界条件。

在这样一类膜中，$A_2-\phi_3$ 与 $A_3+\phi_2$ 有指定的边界值，而
$A_2+\phi_3$ 与 $A_3-\phi_2$ 可以自由涨落。为理解这一结果，回顾
复结构 $J$ 中的全纯坐标是有帮助的：
$2(A_z+i\phi_z)=(A_2-i A_3)+i(\phi_2-i\phi_3)=(A_2+\phi_3)-i(
A_3-\phi_2)$ 与
$2(A_{\bar z}+i\phi_{\bar z})=(A_2+i
A_3)+i(\phi_2+i\phi_3)=(A_2-\phi_3)+i( A_3+\phi_2)$。因此我们的
结果表明，对这类膜，${\cal A}_{\bar
z}=A_{\bar z}+i\phi_{\bar z}$ 被固定，而 ${\cal A}_{ z}=A_{
z}+i\phi_{ z}$ 不受约束。这定义了 $\MH$ 的一个子流形 $W'$，它
显然是复结构 $J$ 下的复子流形。此外，从全纯辛形式
$\Omega_J=(-i/4\pi)\int_C\Tr\, \delta {\cal A}\wedge \delta {\cal
A}$ 的角度看，$W'$ 是 Lagrangian 的。确实，$W'$ 是中间维数的，且
$\Omega_J$ 在 $W'$ 上为零，因为 $W'$ 由
$\delta {\cal A}_{\bar z}=0$ 刻画。于是支撑在 $W'$ 上、带平凡
Chan-Paton 丛的膜就是 $(A,B,A)$ 型膜，符合预期。$W'$ 可刻画为
$\MH$ 中由具有指定全纯结构的平坦 $G_{{\Bbb{C}}}$ 丛组成的
子流形。

对另一类膜，$T'\lambda|=-\lambda|$，超对称要求
$\{\G^{IJ}F_{IJ},T'\}\epsilon=0$。这导出两个条件。第一个是：若
$I,J$ 落在与方向 $1,5,2-7,$ 及 $3+6$ 正交的子空间中，则
$F_{IJ}=0$。特别地，若方向 $I$ 在这个子空间中，则 $F_{0I}=0$，我们
对这样的 $I$ 施加 Dirichlet 边界条件。因此特别地，$A_2+\phi_3$ 与
$A_3-\phi_2$ 满足 Dirichlet 边界条件。$\{\G^{IJ}F_{IJ},T'\}\epsilon=0$
的第二个推论是：$F$ 限制在由方向 $1,5,2-7,3+6$ 生成的子空间上时
必须反自对偶。这导出修正的 Neumann 边界条件（作用于 $A_5,$
$A_{2-7}$, 与 $A_{3+6}$） % 式 (12.29)
\setcounter{equation}{28}
\begin{equation}\label{eq:nono}
\begin{aligned}
F_{15}+F_{2-7,3+6}&=0\\ F_{1,2-7}+F_{3+6,5} & =0\\
F_{1,3+6}+F_{5,2-7}&=0.\\\end{aligned}
\end{equation}

对这种膜，$A_{2+7}$ 与 $A_{3-6}$ 在边界上被固定，而 $A_{2-7}$ 与
$A_{3+6}$ 可以自由涨落。因此从低能看，与 $T'\lambda|=\lambda|$ 的
膜相比，$\CA_z$ 与 $\CA_{\bar z}$ 的角色互换了。于是我们得到支撑
在 $\MH$ 的子流形 $W''$ 上的膜，$W''$ 参数化具有固定反全纯结构
的平坦 $G_{\Bbb{C}}$ 丛。这仍是 $(A,B,A)$ 型膜。

还有另一种更平凡的办法交换 $\CA_z$ 与 $\CA_{\bar z}$ 的角色：把
$\phi_2$ 与 $\phi_3$ 变号，换言之把 $T'$ 换成
$T''=-\Gamma_{15}\Gamma_{2+7}\Gamma_{3-6}$。所以我们实际上有两族
膜，分别支撑在参数化固定全纯结构或反全纯结构的丛的流形上。这两
族的差别在于对远离 $\MH$ 奇点处有质量的模式所加的条件；在 $\MH$
的奇点处它们可能不同。

\par\medskip\emph{$S$-对偶与典范余各向同性膜}

对刚才构造的例子施加 $S$-对偶，可以得到新的膜，它们同样是
$(A,B,A)$ 型。我们知道如何对未破缺超对称的条件施加 $S$-对偶：
$T'\epsilon=\epsilon$ 被换成 $\hat T\epsilon=\epsilon$，其中 % 式 (12.30)
\setcounter{equation}{29}
\begin{equation}\label{eq:nofo}
\hat T ={1+\HG\over \sqrt
2}T'{1-\HG\over \sqrt 2}=-{\Gamma_{15}\over
2}\left((\Gamma_{26}+\Gamma_{37})-\HG(\Gamma_{23}+\Gamma_{67})\right).
\end{equation}

要以这种方式得到新的 $(A,B,A)$ 型膜，我们必须构造出保持满足
$\hat T\epsilon=\Gamma_{2367}\epsilon=\epsilon$ 的超对称的边界
条件。

我们寻找一个满足 $\hat T\lambda|=-\lambda|$ 的超对称膜。经过一段
颇长的计算，从超对称条件 $\{\G^{IJ}F_{IJ},\hat T\}\epsilon=0$ 推出，
玻色子上的边界条件必须是
% 式 (12.31)
\setcounter{equation}{30}
\begin{equation}\label{eq:tomgo}
\begin{aligned}F_{15}+F_{26}+F_{37}& = 0 \\
                    F_{12}+F_{06}-F_{56}& = 0 \\
                    F_{13}+F_{07}-F_{57}& = 0 \\
                    F_{16}+F_{52}-F_{02}& = 0 \\
                    F_{17}+F_{53}-F_{03}& = 0 ,\\\end{aligned}
\end{equation}

同时还有对 $J=4,8,9$ 的 $F_{0J}=0$。

我们在边界上把有质量的场 $A_J,\,J=4,8,9$ 取为零。除此之外，所有场
满足修正的 Neumann 边界条件，在边界上可自由涨落，所以这样得到的
$(A,B,A)$ 型膜是空间充满的膜，其靶空间是整个 $\MH$。

我们以前见过混合的 Neumann-Dirichlet 边界条件，但这里有一个
关键差别。在先前的例子里，如果把远离奇点处不参与低能理论的场
$A_4,A_5,A_8$, 与 $A_9$ 取为零，边界条件就变成纯 Dirichlet 或纯
Neumann，因此在膜的有效二维场论中无关紧要。在当前的例子里却
并非如此。丢掉「有质量」场并选取规范 $A_1=0$ 后，边界条件变为 % 式 (12.32)
\setcounter{equation}{31}
\begin{equation}\label{eq:bdrycc}
\begin{aligned}
D_0\phi_2+\partial_1 A_2&=0,\\ D_0\phi_3+\partial_1 A_3&=0,\\
F_{02}-\partial_1\phi_2&=0,\\ F_{03}-\partial_1\phi_3&=0.\end{aligned}
\end{equation}

当膜上的规范场不是平坦的时候，sigma 模型场就会出现混合的
Dirichlet-Neumann 边界条件；具体地说，边界条件为 % 式 (12.33)
\setcounter{equation}{32}
\begin{equation}\label{eq:zembro}
f_{MN}\partial_0 X^N+g_{MN}\partial_1 X^N=0,
\end{equation}
 其中 $f$ 是靶空间规范场的场强，$g$ 是靶空间度规。与
\eqref{eq:bdrycc} 比较，我们得到 % 式 (12.34)
\setcounter{equation}{33}
\begin{equation}\label{eq:membro}
f={1\over 2\pi}\int_C |d^2z| \Tr\, \left(\delta \phi_z\wedge\delta
A_\bz+\delta\phi_\bz\wedge \delta A_z\right)=\omega_J.
\end{equation}
 这正是靶空间为 $\MH(G,C)$ 的 sigma 模型中典范余各向同性
$(A,B,A)$ 型膜的靶空间规范场。所以这就是我们找到的膜。

\par\medskip\emph{典范余各向同性膜的边界可观测元}

现在我们对典范余各向同性膜的边界可观测元作一个简单的观察。
考虑插入在 $M$ 边界上一点 $P$ 处的局部算子 ${\EUO}$。边界当然
看起来像 $\Bbb{R}\times C$。我们要考虑复结构 $K$ 中 $A$-模型的
拓扑超荷 $Q$ 作用在插入于 $\Bbb{R}\times C$ 中一点 $s\times p$ 的
局部算子上的上同调。$C$ 的度规模 $\{Q,\dots\}$ 无关，故可设 $C$
在 $p$ 附近平坦。

我们可以按局部算子的维数、以及它们在 $p$ 点周围 $C$ 的旋转下
变换的「自旋」来分类。（讨论是局部的，所以 $C$ 的整体结构通常
破坏旋转对称这一点并无妨碍。）扭化理论中的自然维数：$A,\phi,
\psi,\tilde\psi$ 为 1，$\chi,\eta,\tilde\eta$ 为 2。

我们考虑维数为 $n$、自旋为 $n$ 的 $Q$-不变算子。要具有这些量子
数，规范不变算子只能由 $\phi_z$, $\psi_z$, $\tilde\psi_z$ 与协变
导数 $D_z$ 构造。但算子要 $Q$-不变，$D_z$ 就不能出现，因为即使
在边界上 $A_z$ 也不是 $Q$-不变的。边界条件 $\hat T\psi|=-\psi|$
意味着
% 式 (12.35)
\setcounter{equation}{34}
\begin{equation}\label{eq:bdrry}
\begin{aligned}\psi_2|&=-i\tilde\psi_3|\\
\psi_3|&=i\tilde\psi_2|.\\\end{aligned}
\end{equation}
 由 $\psi_z=(\psi_2-i\psi_3)/2$ 得到
$\psi_z|=\tilde\psi_z|$。因此维数为 $n$、自旋为 $n$ 的 $Q$-不变
边界可观测元只是 $\phi_z$ 与 $\psi_z$ 的函数，不含任何导数。

场 $\phi_z$ 与 $\psi_z$ 的鬼数分别为 $\EUK=0$ 与 $\EUK=1$。所以
这类算子若还满足 $\EUK=0$，就只是 $\phi_z$ 的规范不变函数，例如
${\EUO}=\Tr\,\phi_z^n$（某个 $n$）。$\EUK=1$ 的典型算子是
${\EUO}'={\rm Tr}\,\phi_z^{n-1}\psi_z$。

经典地，边界条件保证在边界上 $[Q,\phi_z\,]=\{Q,\psi_z\,\}=0$，所以
刚才描述的所有算子 ${\EUO}$ 与 ${\EUO}'$ 都是 $Q$ 上同调中的非零
元。量子力学上会怎样？$Q$ 的作用会不会有量子修正，使得精确的
量子公式为 $[Q,{\EUO}]=\epsilon{\EUO}'$（某个 $\epsilon$）？若如此，
在量子层面 ${\EUO}$ 与 ${\EUO}'$ 就会配对并从同调中消失。

可以用时间反演对称性论证这不会发生。我们定义一个反定向的对称性
${\EUT}$，它把时间坐标 $x^0$ 变号，并平凡地作用在 $M$ 的其它坐标
上。考虑到 ${\EUT}$ 作用在 $\phi$ 上会带一个负号（见式
\eqref{eq:yoro}），${\EUT}$ 是定义典范余各向同性膜的边界条件的
对称性。${\EUT}$ 与复结构 $K$ 中 $A$-模型的拓扑超荷对易，因为
$t=1$ 是 ${\EUT}$ 的不动点（见式 \eqref{eq:zongor}）。又由
\eqref{eq:yoro}，$\phi_z$ 在 ${\EUT}$ 下为奇，但 \eqref{eq:zoro}
加上边界条件 $\psi_z|=\tilde\psi_z|$ 意味着 $\psi_z|$ 在 ${\EUT}$
下为偶。所以算子 ${\EUO}$ 与 ${\EUO}'$ 不可能在 $Q$ 的作用下配对，
必须在上同调中存活。更一般地，考虑维数与自旋为 $n$、鬼数为 $k$
的局部边界算子，它在经典层面是 $Q$-不变的。这样的算子由规范不变
函数 $F(\phi_z,\psi_z)$ 决定，它在 $\phi_z$ 中的次数为 $n-k$，在
$\psi_z$ 中的次数为 $k$。它在 ${\EUT}$ 下按 $(-1)^{n-k}$ 变换。
若 $Q$ 在量子层面要非平凡地作用在这样的算子上，它必须保持 $n$
不变并把 $k$ 增加 1。这会使 ${\EUT}$ 的本征值反号，因此不可能。
于是 $Q$ 在维数与自旋为 $n$ 的局部算子空间中的上同调与经典结果
一致，正是规范不变函数 $F(\phi_z,\psi_z)$ 构成的空间。

\begin{yizhu}[易错点]
$S$-对偶把 $T$ 换成 $T'$ 时，投影算子的次序不可交换：
$T'={1+\HG\over\sqrt 2}\,T\,{1-\HG\over\sqrt 2}$（式 \eqref{eq:uvuf}），
两个因子的 $\pm$ 不能随手写成同一个号；上式成立的前提是
$T\epsilon=\epsilon$，此时 $T'$ 才作用在 $\epsilon'$ 上满足
$T'\epsilon'=\epsilon'$。由此 $T=\Gamma_0\Gamma_4$（$(B,B,B)$ 型膜的
标准选择）给出 $T'=\Gamma_1\Gamma_2\Gamma_3\Gamma_4$（式
\eqref{eq:polyo}），后者的四个方向全是「$A$ 型」方向，这正是
$(B,A,A)$ 型膜出现的原因：$S$-对偶把原来的两个 Neumann 方向
$0,4$ 换成了 $1,2,3,4$ 中的三个外加一个共同的 $\Gamma_4$。
\end{yizhu}

\begin{yizhu}[易错点]
本片有两类容易混淆的边界数据。其一是场强与势：$F_{0I}|=0$ 是
场强分量的 Dirichlet 型条件，而 $A_I\vert=0$ 是势本身取零，原文
对 $I=5,8,9$ 直接令 $A_I\vert=0$，对 $I=6,7$ 则只固定
$\phi_z={1\over 2}(\phi_2-i\phi_3)$ 的特征多项式（而非把
$A_6,A_7$ 取零），所以 $A_6,A_7$ 仍是「Dirichlet 型」的：它们的
边界值被指定，但不是零。其二是（12.25）的修正 Neumann 条件给出
的是 $F_{12}\vert-F_{34}\vert=0$ 这类分量之差为零，而不是单个
$F_{ij}$ 消失；这正是缠绕 Hitchin 纤维所需的条件。
\end{yizhu}

\begin{yizhu}[含糊点拨]
原文说 $D_\bz\phi_z=0$ 是「复」Hitchin 方程，引号不可省：真正的
Hitchin 方程是 $F-\phi\wedge\phi=0$ 与 $D\phi=D\star\phi=0$，而这里
只用到复结构 $J$ 中方程的一个全纯半支，且（12.26）只是两个实
方程。按：这个「半支」性质正是该膜支撑在 Hitchin 纤维上、而不
在 $\MH$ 内部的原因；它与第~\ref{sec:xgoto}~节后面关于 $W'$、$W''$
是 Lagrange 子流形的论断相呼应。
\end{yizhu}

\begin{yizhu}[含糊点拨]
（12.28）引入的 $\Gamma_{\pm I\pm J}$ 与随后大量出现的指标记法
（如 $F_{2-7,3+6}$、$A_{2-7}$）中，$2-7$ 表示方向 $2$ 与方向 $7$
的组合 $\Gamma_2-\Gamma_7$，是缩写而非减法；同理 $3+6$ 表示
$\Gamma_3+\Gamma_6$。读（12.29）时必须按这一约定展开，否则会把
$2-7$ 误读成第七维以外的某个负指标。
\end{yizhu}

\begin{yizhu}[物理图像]
（12.34）给出的是膜的「靶空间规范场」$f$：把（12.33）的一般
混合 Dirichlet-Neumann 条件与（12.32）对比，可见在靶空间为
$\MH(G,C)$ 的 sigma 模型中，典范余各向同性 $(A,B,A)$ 型膜对应的
正是 $f=\omega_J$（复结构 $J$ 下的 Kähler 形式）。按：这就是
「典范余各向同性」一名的来源——膜的支撑流形是空间充满的，而
其上的 B-场/规范场取为辛形式本身；它与第~\ref{sec:zapendix}~节中
用 $A=\mathrm{Hom}(A_{c.c.},A)$ 定义典范余各向同性膜的说法一致。
\end{yizhu}

\begin{yizhu}[物理图像]
最后一段的时间反演论证值得记住其逻辑顺序：${\EUT}$ 是定义
典范余各向同性膜的边界条件的对称性（故可用来约束边界可观测
元），它与复结构 $K$ 中 $A$-模型的拓扑超荷对易（因为 $t=1$ 是
不动点），而 $\phi_z$ 在 ${\EUT}$ 下为奇、$\psi_z$ 为偶，于是
$F(\phi_z,\psi_z)$ 按 $(-1)^{n-k}$ 变换。任何量子修正若要非平凡
地作用，必须把 $k$ 加一，从而翻转该本征值，故不可能发生。按：
这一论证说明边界可观测元的上同调是「经典的」，不含量子修正，
因此 $\Tr\phi_z^n$ 这类算子可以当作 $(A,B,A)$ 型膜的真正边界
自由度。
\end{yizhu}

\begin{yizhu}[原文疑误]
两处照录并说明。（1）式 \eqref{eq:membro} 右端把 $C$ 上的积分
直接等同于辛形式 $\omega_J$，而左端的 $f$ 在（12.33）中是靶空间
规范场（二维场的场强），原文未交代 $2\pi$ 因子与量纲的来历，
按：这里应理解为一个定义式的写法，$f:=\omega_J$ 才是实质内容。
（2）"archetypical" 一词通常作 "archetypal"；又本片与上一片
交替使用宏 \verb|\eps| 与 \verb|\epsilon| 表示同一超对称参数，
原文如此，译文照录。
\end{yizhu}
